\section{Spherical Sector Geometry}

\subsection{Centering and normalization}

Let $\mathbf{1}$ denote the all-ones vector and define
\[
M_c = M - \frac{14}{30}\,\mathbf{1}\mathbf{1}^{\mathsf T}.
\]

Then the rows of $M_c$ lie in the $14$-dimensional subspace orthogonal to
$\mathbf{1}$. Normalize by setting
\[
\widetilde{M}_v = \frac{(M_c)_v}{\|(M_c)_v\|}.
\]

\subsection{Three-angle structure}

\begin{theorem}[Exact three-angle structure]
The vectors $\{\widetilde{M}_v\}$ form a spherical configuration in
$\mathbb{R}^{14}$ with exactly three distinct inner products.
More precisely, if $u \neq v$, then
\[
\langle \widetilde{M}_u, \widetilde{M}_v \rangle
=
\begin{cases}
\frac{37}{112} & \text{if } d(u,v)=1, \\[4pt]
-\frac{23}{112} & \text{if } d(u,v)=2, \\[4pt]
-\frac{38}{112} & \text{if } d(u,v)=3,
\end{cases}
\]
where $d(u,v)$ denotes graph distance in $\Gfifteen$.
\end{theorem}

\begin{proof}
Let
\[
B := M M^{\mathsf T} = A^3 + 2A^2 + 2I,
\]
where $A$ is the adjacency matrix of
\[
\Gfifteen \cong L(\mathrm{Petersen}).
\]

Since $\Gfifteen$ is distance-regular with intersection array
\[
\{4,2,1;1,1,4\},
\]
its distance matrices $D_0,D_1,D_2,D_3$ satisfy
\[
A^2 = 4D_0 + D_1 + D_2,
\qquad
A D_2 = 2D_1 + 2D_2 + 4D_3.
\]
It follows that
\[
A^3 = 4D_0 + 7D_1 + 3D_2 + 4D_3.
\]
Substituting into $B$ gives
\[
B = 14D_0 + 9D_1 + 5D_2 + 4D_3.
\]

Since $A\mathbf{1}=4\mathbf{1}$, we have
\[
B\mathbf{1} = (4^3 + 2\cdot4^2 + 2)\mathbf{1} = 98\,\mathbf{1},
\]
so every row sum of $B$ equals $98$. Let
\[
H = I - \frac{1}{15}J.
\]
Then the centered Gram matrix is
\[
HBH = B - \frac{98}{15}J.
\]
Using
\[
J = D_0 + D_1 + D_2 + D_3,
\]
we obtain
\[
HBH
=
\frac{112}{15}D_0
+\frac{37}{15}D_1
-\frac{23}{15}D_2
-\frac{38}{15}D_3.
\]

The common squared norm is $\frac{112}{15}$, so after normalization the
Gram matrix becomes
\[
G
=
D_0
+\frac{37}{112}D_1
-\frac{23}{112}D_2
-\frac{38}{112}D_3,
\]
which yields the stated inner products.
\end{proof}

\subsection{Geometric interpretation}

The embedding
\[
v \longmapsto \widetilde{M}_v
\]
realizes $\Gfifteen$ as a rigid spherical three-angle configuration in
$\mathbb{R}^{14}$.

The three inner products correspond exactly to graph distances
$1,2,3$, showing that this geometry is completely determined by the
quadratic overlap structure of the transport sectors.

\subsection{Geometric realization of the cubic representation}

This embedding is the geometric realization of the cubic Thalean
representation. The polynomial identity
\[
M M^{\mathsf T} = p(A)
\]
translates directly into metric relations in the embedding, making the
sector geometry canonical at the quadratic level.

As discussed in Section~11, this spherical geometry depends only on the
sector incidence structure and does not by itself detect the cocycle.
